% =====================================================================
\section{Cotta \#2 --- Witbier}
% =====================================================================

\begin{tabularx}{\linewidth}{@{}lX@{}}
\toprule
Data cotta & 25/10/2026 \\
Stile & Witbier (BJCP 2021, 24A, p.~49) \\
Parametri di stile & OG 1,044--1,052; FG 1,008--1,012; IBU 8--20; SRM 2--4; ABV 4,5--5,5\,\% \\
Volume in fermentatore & \SI{17}{\liter} ($4 + 4 + \SI{9}{\liter}$) \\
Volume di riferimento ricetta & \SI{23}{\liter} (ricetta di partenza, scalata a \SI{17}{\liter}) \\
Efficienza & $\varepsilon = 0{,}78$ usata in pianificazione (cotta \#1); \textbf{misurata: $\varepsilon = 0{,}63$} \\
Costo ingredienti & 32,06\,\texteuro\ + zucchero di priming \textcolor{red}{Da definire} (1,89\,\texteuro/L su \SI{17}{\liter}, escluso lo zucchero; esclusa la bottiglia; incluso mezzo barattolo di sanificante) \\
Carbonazione obiettivo & $v_t = \SI{3,0}{\vol}$ (BJCP 24A: \emph{high carbonation}, senza valore numerico; scelto prudente per bottiglie da \SI{33}{\centi\liter}) \\
\bottomrule
\end{tabularx}

\subsection*{Obiettivi}
\begin{tabularx}{\linewidth}{@{}lccX@{}}
\toprule
Grandezza & Obiettivo & Misurato & Note \\
\midrule
$V_{pb}$ [\si{\liter}] & 20 & 18 & a freddo. $V_{fer} + V_l$: $V_l = \SI{1}{\liter}$ misurato ($18 - 17$), contro i \SI{3}{\liter} di cotta \#1 (bollitore inclinato) \\
$OG$ & 1,047 & 1,042 & misurato a \SI{19}{\celsius} (lettura effettiva, offset del nuovo densimetro già applicato); $GU_{pb} = 42$, sotto l'obiettivo. Pre-bollitura: 1,035 a \SI{22}{\liter} caldi (lettura 1,032 a \SI{32}{\celsius}, più correzione $+3{,}1$). Range dello stile 1,044--1,052 \\
$P_{tot}$ [pt$\cdot$L] & 940 & 756 & misurato: $42 \cdot 18$ (obiettivo $47 \cdot 20$) \\
$FG$ & 1,012 & \textcolor{red}{Da definire} & $A_{app} = 0{,}75$ (limite superiore dichiarato del lievito); con $OG = 1{,}042$ la FG attesa è $\approx 1{,}011$; range dello stile 1,008--1,012 \\
$IBU$ & $\approx 17$ & $\approx 19$ & ricetta: $\approx 17$ a \SI{20}{\liter}. Misurato: stima con \SI{18}{\liter} (Bobek $14{,}0$ + Saaz $4{,}7$), utilizzi teorici, non misura diretta; dentro il range dello stile (8--20) \\
$ABV$ [\si{\percent}] & $\approx 4{,}6$ & \textcolor{red}{Da definire} & prima del priming; $ABV = 131{,}25\,(OG-FG)$. Con $OG = 1{,}042$ e FG attesa $1{,}011$: $\approx 4{,}1$\,\% prima del priming, sotto il minimo di stile. Dopo il priming $\approx 4{,}7$\,\% \\
\bottomrule
\end{tabularx}

\noindent \textbf{Nota:} con $A_{app} = 0{,}75$ il FG cade sul limite superiore del range. Se la FG reale risulta più alta, l'ABV scende sotto 4,5\,\%: in quel caso la prossima cotta va alzata di OG. Con l'OG reale ($1{,}042$) l'ABV atteso prima del priming è già sotto 4,5\,\%: la FG va misurata.

% ---------------------------------------------------------------------
\subsection*{Ricetta di partenza e modifiche}

\begin{tabularx}{\linewidth}{@{}lXXX@{}}
\toprule
Voce & Partenza (\SI{23}{\liter}) & Cotta \#2 (\SI{17}{\liter}) & Motivo \\
\midrule
OG & 1,045 & 1,047 & ABV $\ge 4{,}5\,\%$ con FG $\le 1{,}012$ \\
Frumento non maltato & \SI{0,7}{\kilo\gram} (15\,\%) & \SI{1,26}{\kilo\gram} (32\,\%) & BJCP 24A: 30--60\,\% \\
Frumento maltato (Château Wheat Blanc) & \SI{1,2}{\kilo\gram} & \SI{0,95}{\kilo\gram} & bilancio sull'OG obiettivo, con 1,26 kg di frumento non maltato \\
Château Pilsner & \SI{2,5}{\kilo\gram} & \SI{1,49}{\kilo\gram} & bilancio sull'OG obiettivo \\
Fiocchi d'avena & \SI{0,3}{\kilo\gram} & \SI{0,27}{\kilo\gram} & circa 7\,\% della miscela (stile: fino al 5--10\,\%) \\
Luppolo amaro & Bobek \SI{25}{\gram}, 4,5\,\% & Bobek \SI{20}{\gram}, 4,2\,\% & IBU obiettivo 16: margine sotto il limite superiore dello stile (20) \\
Luppolo aroma & Saaz \SI{25}{\gram}, 3,2\,\% & Saaz \SI{16}{\gram}, 3,5\,\% & luppolo speziato \emph{low to none} nello stile \\
Arancia amara & \SI{20}{\gram} & \SI{15}{\gram} & scala a \SI{17}{\liter} (invariata rispetto alla versione a 19 L) \\
Coriandolo & \SI{10}{\gram} & \SI{7}{\gram} & scala a \SI{17}{\liter}; lo stile vuole spezia non esagerata \\
\bottomrule
\end{tabularx}

% ---------------------------------------------------------------------
\subsection*{Ricetta}

\subsubsection*{Malto}
\begin{tabularx}{\linewidth}{@{}Xrr@{}}
\toprule
Malto & Massa [\si{\kilo\gram}] & \% \\
\midrule
Château Pilsner & 1,49 & 37,5 \\
Château Wheat Blanc (frumento maltato) & 0,95 & 23,9 \\
Torrified Wheat (frumento non maltato) & 1,26 & 31,8 \\
Fiocchi d'avena & 0,27 & 6,8 \\
\midrule
Totale & 3,97 & 100,0 \\
\bottomrule
\end{tabularx}

\noindent Calcolo: $m_i = P_i / (10\, e_i\, \varepsilon)$, con $P_{tot} = 940$ pt$\cdot$L ripartiti in Pilsner ($P = 349$), frumento maltato ($P = 237$), frumento non maltato ($P = 295$) e avena ($P = 59$). Estratti potenziali $e_i$ dalla tabella del formulario: 30 (Pilsner), 32 (frumento maltato, centro del range 31--33), 30 (frumento non maltato), 28 (avena).

\subsubsection*{Luppolo}
\begin{tabularx}{\linewidth}{@{}Xccccr@{}}
\toprule
Luppolo & Forma & $a$ [\si{\percent}] & Massa [\si{\gram}] & Aggiunta & $IBU_i$ \\
\midrule
Bobek & pellet & 4,2 & 20 & \SI{60}{\minute} & 12,6 \\
Saaz & pellet & 3,5 & 16 & \SI{15}{\minute} & 4,2 \\
\bottomrule
\end{tabularx}

\noindent Masse calcolate con rapporto $r = m_{h,1}/m_{h,2} = 1{,}2$, $IBU = 16$, $C_g = 1$ (OG $< 1{,}050$), $V_{pb} = \SI{20}{\liter}$, utilizzi $U=0{,}30$ (60 min) e $U=0{,}15$ (15 min) in pellet.

\subsubsection*{Altri ingredienti}
\begin{tabularx}{\linewidth}{@{}Xcr@{}}
\toprule
Ingrediente & Aggiunta & Quantità [\si{\gram}] \\
\midrule
Arancia amara (buccia essiccata) & \SI{15}{\minute} & 15 \\
Coriandolo (semi) & \SI{5}{\minute} & 7 \\
\bottomrule
\end{tabularx}

\noindent Quantità scalate dalla ricetta di partenza ($20\,\mathrm{g}$ e $10\,\mathrm{g}$ per \SI{23}{\liter}).

\subsubsection*{Lievito}
\begin{tabularx}{\linewidth}{@{}lX@{}}
\toprule
Ceppo & Mangrove Jack's Belgian Wit M21 \\
Quantità & \SI{10}{\gram} $\;\to\; q_y = \SI{0,59}{\gram\per\liter}$ \\
Attenuazione dichiarata & 70--75\,\% (usata $75$\,\% per il calcolo di FG, vedi nota sopra) \\
Range temperatura & 18--\SI{25}{\celsius} \\
Dose massima dichiarata & fino a \SI{23}{\liter} di mosto \\
\bottomrule
\end{tabularx}

\subsubsection*{Acqua}
\begin{tabularx}{\linewidth}{@{}lccc@{}}
\toprule
 & Rocchetta [\si{\liter}] & Lauretana [\si{\liter}] & Totale [\si{\liter}] \\
\midrule
Mash   & 8,5 & 5,5 & 14,0 \\
Sparge & 8,0 & 5,5 & 13,5 \\
\midrule
Totale & 16,5 & 11,0 & 27,5 \\
\bottomrule
\end{tabularx}

\medskip
\noindent Miscela $\approx$ 60\,\% Rocchetta / 40\,\% Lauretana (come cotta \#1). Profilo ionico uguale a cotta \#1: stessa miscela complessiva (60\,\% / 40\,\% su \SI{27,5}{\liter}), con mash e sparge che differiscono di meno di 2 punti percentuali. Sparge acidificata a pH $\approx 5$.

% ---------------------------------------------------------------------
\subsection*{Processo}

\subsubsection*{Volumi}
\begin{tabularx}{\linewidth}{@{}lX@{}}
\toprule
Volume totale acqua & $V_{tot} = V_{fer} + V_l + V_{evp} + V_{abs} = 17 + 3 + 3 + 4{,}5 \approx \SI{27,5}{\liter}$ \\
Assorbimento trebbie & $\eta = \SI{1,125}{\liter\per\kilo\gram} \Rightarrow V_{abs} \approx \SI{4,5}{\liter}$ (su $\SI{3,97}{\kilo\gram}$) \\
Evaporazione & $\sigma = \SI{3}{\liter\per\hour}$ per \SI{60}{\minute} \\
Rapporto di mash $R$ & $\SI{3,5}{\liter\per\kilo\gram} \Rightarrow V_m \approx \SI{13,9}{\liter}$ (pianificato; effettivo \SI{14,0}{\liter}) \\
Volume sparge & $V_{sp} = V_{tot} - V_m \approx \SI{13,6}{\liter}$ (pianificato; effettivo \SI{13,5}{\liter}) \\
Bilancio misurato & Con contrazione a caldo (fattore $0{,}962$): $V_{pre} = \SI{22}{\liter} \cdot 0{,}962 \approx \SI{21,2}{\liter}$ a freddo. $\sigma = \SI{21,2}{\liter} - \SI{18}{\liter} = \SI{3,2}{\liter\per\hour}$, coerente con $\SI{3}{\liter\per\hour}$ della cotta \#1. $V_{abs} = \SI{27,5}{\liter} - \SI{21,2}{\liter} = \SI{6,3}{\liter} \Rightarrow \eta = \SI{6,3}{\liter} / \SI{3,97}{\kilo\gram} \approx \SI{1,59}{\liter\per\kilo\gram}$. Somma $V_{ev}+V_{abs} = \SI{9,5}{\liter}$ \\
\bottomrule
\end{tabularx}

\subsubsection*{Mash e sparge}
\begin{tabularx}{\linewidth}{@{}Xcc@{}}
\toprule
Step & Temperatura & Durata \\
\midrule
Protein rest & \SI{52}{\celsius} & \SI{15}{\minute} \\
Mash & \SI{62}{\celsius} & \SI{30}{\minute} \\
Mash & \SI{65}{\celsius} & \SI{30}{\minute} \\
Mash & \SI{72}{\celsius} & \SI{10}{\minute} \\
Mash-out & \SI{78}{\celsius} & \SI{10}{\minute} \\
\bottomrule
\end{tabularx}

\subsubsection*{Bollitura}
\begin{tabularx}{\linewidth}{@{}lX@{}}
\toprule
Durata & \SI{60}{\minute} \\
Volume pre-bollitura & \SI{22}{\liter} a caldo (misurato, prima dell'accensione) \\
Volume post-bollitura & \SI{18}{\liter} a freddo (misurato; obiettivo \SI{20}{\liter}) \\
Aggiunte & \SI{0}{\minute}: Bobek \SI{20}{\gram} (60 min) \\
 & \SI{45}{\minute} dopo l'inizio: Saaz \SI{16}{\gram} + arancia amara \SI{15}{\gram} (15 min alla fine) \\
 & \SI{55}{\minute} dopo l'inizio: coriandolo \SI{7}{\gram} (5 min alla fine) \\
\bottomrule
\end{tabularx}

\subsubsection*{Fermentazione}
\begin{tabularx}{\linewidth}{@{}lX@{}}
\toprule
Fase 1 & \SI{19}{\celsius} per 7 gg (fermentazione principale; range effettivo da termostato \SI{18}{}--\SI{19}{\celsius}) \\
Fase 2 & \SI{22}{\celsius} per 7 gg (fine fermentazione, riassorbimento del diacetile; range termostato \SIrange{21}{22}{\celsius}) \\
Temperatura massima $T_f^{max}$ & \SI{22}{\celsius} (range lievito \SIrange{18}{25}{\celsius}) \\
Durata totale & 14 gg, poi imbottigliamento \\
Misurato & \textcolor{red}{Da definire} \\
\bottomrule
\end{tabularx}

\subsubsection*{Priming}
\begin{tabularx}{\linewidth}{@{}lX@{}}
\toprule
Volume CO$_2$ obiettivo $v_t$ & \SI{3,0}{\vol} \\
CO$_2$ residua $v_r(T_f^{max} = \SI{22}{\celsius})$ & \SI{0,80}{\vol} (tabella del formulario) \\
Zucchero & $c_s = 3{,}86\,(3{,}0 - 0{,}80) \approx \SI{8,5}{\gram\per\liter}$ (saccarosio) \\
$\Delta ABV_p = 0{,}0648\, c_s$ & $\approx 0{,}6$\,\% v/v (ABV finale $\approx 4{,}1 + 0{,}6 \approx 4{,}7$\,\%) \\
Volume da imbottigliare $V_b$ & \textcolor{red}{Da definire} [\si{\liter}] \\
Massa di zucchero $m_s = c_s V_b$ & \textcolor{red}{Da definire} [\si{\gram}] \\
\bottomrule
\end{tabularx}

% ---------------------------------------------------------------------
\subsection*{Problemi e note per le prossime cotte}
\begin{itemize}
  \item \textbf{Bolle dal coperchio.} Ci siamo allontanati durante la bollitura, che era ancora in corso. Al ritorno uscivano bolle dal coperchio: è uscito un po' di mosto, perdita piccola e non misurata. La differenza $\SI{21,2}{\liter} - \SI{18}{\liter} = \SI{3,2}{\liter}$ include anche la perdita, quindi l'evaporazione reale è più bassa. L'efficienza non cambia, perché il volume post-bollitura è stato misurato dopo. Non lasciare il bollitore incustodito durante la bollitura; se serve allontanarsi, appoggiare il coperchio in modo che sfiati.
  \item \textbf{Efficienza misurata $\varepsilon = 0{,}63$ contro $0{,}78$.} Le stime di estratto per frumento e avena sono corrette, quindi il mosto conferma il volume misurato: $1{,}035 \cdot \SI{22}{\liter}$ caldi corrispondono a circa \SI{21,2}{\liter} freddi ($0{,}962$), cioè circa $741$ pt$\cdot$L, contro $756$ pt$\cdot$L post-bollitura. Le masse di malto sono corrette, quindi il deficit sta nell'ammostamento, non nelle misure di fine bollitura. Cause da verificare, in quest'ordine: (a) pH di ammostamento: misurato con cartine sia per il mash sia per lo sparge, tra 5 e 6; l'acqua ha $\mathrm{HCO_3^-} = 113$\,mg/L e il mash non è acidificato; (b) macinatura grossa; (c) temperatura o durata delle pause.
  \item \textbf{Sensibilità al volume.} $\varepsilon$ dipende direttamente dal volume post-bollitura: con \SI{20}{\liter} invece di \SI{18}{\liter} l'efficienza sarebbe circa $0{,}70$. Il volume post-bollitura è stato misurato a freddo, alla stessa temperatura del travaso in fermentatore.
  \item \textbf{Densimetro.} Il nuovo strumento legge $1{,}001$ in acqua distillata: errore $+0{,}001$, quindi $\delta = -0{,}001$ da sommare alle letture. Il valore $\delta = +0{,}005$ del vecchio strumento non si applica più.
  \item \textbf{ABV.} Con OG $1{,}042$ e FG attesa $1{,}011$, l'ABV prima del priming è circa $4{,}1$\,\%, sotto il minimo di stile. Per la prossima cotta alzare l'OG.
\end{itemize}
